% ATLAS: what needs your hands now -- the run-book of the jobs that only a person can do (jobs A to E).
% Written by hand; compile with pdflatex. The look (fonts, colours, grey command blocks, tick boxes, note boxes) is the
% one of the Station PC guide. Facts are the team lead's, as of 19:50 on 3 October 2026.
\documentclass[10pt]{article}
\usepackage[a4paper,top=13mm,bottom=17mm,left=16mm,right=16mm,footskip=9mm]{geometry}
\usepackage[T1]{fontenc}
\usepackage[utf8]{inputenc}
\input{glyphtounicode}\pdfgentounicode=1
\usepackage{helvet}
\renewcommand{\familydefault}{\sfdefault}
\usepackage[varqu,varl]{zi4}
\usepackage[table]{xcolor}
\usepackage{array,booktabs}
\usepackage{enumitem}
\usepackage{needspace}
\usepackage{listings}

\pdfinfo{/Title (ATLAS: what needs your hands now)}

% ---------- colours (the palette of the other ATLAS handouts) ----------
\definecolor{ink}{HTML}{1F2933}
\definecolor{mute}{HTML}{56616B}
\definecolor{laneD}{HTML}{E3EDFC}
\definecolor{accC}{HTML}{9A4A00}
\definecolor{accD}{HTML}{1B55B0}
\definecolor{cmdbg}{HTML}{F1F3F5}
\definecolor{trapbg}{HTML}{FDF3E7}
\definecolor{laterbg}{HTML}{EEF2F6}

% ---------- footer ----------
\makeatletter
\def\ps@atlas{%
  \def\@oddhead{}\def\@evenhead{}%
  \def\@oddfoot{\parbox{\textwidth}{\footnotesize\color{mute}ATLAS: what needs your hands now\hfill 3 October 2026\hspace{6mm}Page \thepage}}%
  \let\@evenfoot\@oddfoot}
\makeatother
\pagestyle{atlas}

\setlength{\parindent}{0pt}
\setlength{\parskip}{0pt}
\color{ink}

% ---------- commands to type, and what they print ----------
% Commands are set so that copying them out of the PDF gives exactly what is printed:
% plain ASCII quotes and hyphens, one command per line, no prompt characters, no wrapped lines.
\lstdefinestyle{cmd}{basicstyle=\ttfamily\fontsize{8.35}{11.3}\selectfont\color{ink},columns=fullflexible,
  keepspaces=true,upquote=true,showstringspaces=false,backgroundcolor=\color{cmdbg},frame=none,
  aboveskip=3.5pt,belowskip=3.5pt,xleftmargin=6pt,framexleftmargin=6pt,breaklines=false}
\lstdefinestyle{inl}{basicstyle=\ttfamily\fontsize{8.8}{11}\selectfont,columns=fullflexible,keepspaces=true,
  upquote=true,showstringspaces=false,breaklines=true}
\lstnewenvironment{cmd}{\lstset{style=cmd}}{}
\newcommand{\cc}{\lstinline[style=inl]}
\newcommand{\code}[1]{{\ttfamily\fontsize{8.8}{11}\selectfont #1}}
\newcommand{\see}{\par\vspace{0.6mm}{\footnotesize\color{mute}You should see:}}

% ---------- jobs and their steps ----------
\newlength{\cbw}      % width of a tick box and its gap; measured after \begin{document}
\newcommand{\jobletter}{}
% \Job{letter}{title}{one line: why}   the blue band of the Station guide's parts, with the letter of the job
\newcommand{\Job}[3]{%
  \par\vspace{6mm}\needspace{52mm}\gdef\jobletter{#1}%
  {\setlength{\fboxsep}{0pt}\colorbox{laneD}{\parbox{\linewidth}{\vspace{2.2mm}\hspace{3mm}%
    {\fontsize{17}{19}\selectfont\bfseries\color{accD}#1}\hspace{3mm}%
    {\fontsize{13.5}{16}\selectfont\bfseries\color{ink}#2}\vspace{2.2mm}}}}\par\nopagebreak
  \vspace{1.8mm}{\small\color{mute}#3\par}\nopagebreak\vspace{1.4mm}}
% the steps of a job: a tick box, then the job letter and a number (A1, A2, ...)
\newlist{jsteps}{enumerate}{1}
\setlist[jsteps]{label=\cb\textbf{\jobletter\arabic*},leftmargin=12.5mm,labelwidth=10.4mm,labelsep=2.1mm,align=left,itemsep=2.8mm,topsep=1.5mm,parsep=0.8mm}
\newcommand{\act}[1]{\textbf{#1}\enskip}
% a step keeps its first lines and its first command on one page
\newcommand{\step}{\needspace{21mm}\item}
% ticked sub-steps inside one step
\newlist{ticks}{itemize}{1}
\setlist[ticks]{label=\cb,leftmargin=\cbw,labelsep=0mm,align=left,itemsep=1.6mm,topsep=1.2mm,parsep=0.6mm}

% ---------- boxes: a trap, and a thing that waits ----------
\newsavebox{\notebox}
\newenvironment{trap}[1][If it fails]{\par\vspace{1.4mm}\begin{lrbox}{\notebox}%
  \begin{minipage}{\dimexpr\linewidth-5mm\relax}\small\raggedright{\bfseries\color{accC}#1}\enskip}
  {\end{minipage}\end{lrbox}{\setlength{\fboxsep}{2.5mm}\colorbox{trapbg}{\usebox{\notebox}}}\par\vspace{1.4mm}}
\newenvironment{later}[1][Waits for code]{\par\vspace{1.4mm}\begin{lrbox}{\notebox}%
  \begin{minipage}{\dimexpr\linewidth-5mm\relax}\small\raggedright{\bfseries\color{accD}#1}\enskip}
  {\end{minipage}\end{lrbox}{\setlength{\fboxsep}{2.5mm}\colorbox{laterbg}{\usebox{\notebox}}}\par\vspace{1.4mm}}

% ---------- small furniture ----------
\newcommand{\cb}{\raisebox{-0.25ex}{\setlength{\fboxsep}{0pt}\setlength{\fboxrule}{0.7pt}\fcolorbox{mute}{white}{\rule{0pt}{1.75ex}\rule{1.75ex}{0pt}}}\hspace{2mm}}
\newcommand{\h}[1]{\par\vspace{4mm}{\fontsize{11.5}{14}\selectfont\bfseries\color{ink}#1}\par\vspace{1.4mm}}
\newcommand{\m}{\,\textrm{\guilsinglright}\,}
% one row of the overview: a tick box, then the letter of the job
\newcommand{\jobrow}[1]{{\normalsize\cb}\textbf{\color{accD}#1}}

\begin{document}
\raggedright
\settowidth{\cbw}{\cb}

% =====================================================================
% PAGE 1: the title, the state of the Mac, and the five jobs
% =====================================================================
{\fontsize{26}{30}\selectfont\bfseries ATLAS: what needs your hands now\par}
\vspace{2.5mm}
{\large\color{mute}What only a person can do right now, step by step.\\
Everything else is being written by the agent sessions.\par}
\vspace{2.5mm}
{\small\color{mute}State at 19:50 on 3 October 2026. The Station PC is set up. On the Mac: Xcode 26.1 is selected; no Apple account
is signed in to Xcode; there is no Apple~Development~certificate; no phone has shown as ready yet.\par}

\h{The five jobs}
{\small
\renewcommand{\arraystretch}{1.45}
\begin{tabular}{@{}l@{\hspace{3mm}}>{\bfseries\raggedright\arraybackslash}p{50mm}@{\hspace{3mm}}>{\raggedright\arraybackslash}p{108mm}@{}}
\toprule
 & Job & What it is, and what it blocks\\
\midrule
\jobrow{A} & Sign in to Xcode & Blocks every phone build.\\
\jobrow{B} & Prepare the phones & Blocks the first build and all testing on a phone.\\
\jobrow{C} & First build to the phones & The first proof that signing and the phones work.\\
\jobrow{D} & Decisions only you can make & Four decisions. None is urgent this hour.\\
\jobrow{E} & The agent sessions & Neither branch has reached GitHub yet.\\
\bottomrule
\end{tabular}}

\begin{later}[How to read it.]
Tick a box when a step is done. Grey blocks are copied out of this PDF and used exactly as printed. Unless a step says
otherwise, they go into Terminal on the Mac. A small line that says \emph{You should see:} tells you what a good result
looks like.
\end{later}

% =====================================================================
\Job{A}{Sign in to Xcode}{Blocks every phone build.}
{\small\color{mute}Already done: the command-line tools point at Xcode, and \code{xcodebuild -version} prints \code{Xcode 26.1}.\par}
\vspace{0.5mm}
\begin{jsteps}
\step \act{Open Xcode\m Settings (Cmd + ,)\m Accounts.} Click +\m Apple Account and sign in. A free account gives a
      Personal Team.
      \see\ the account in the list, with ``(Personal Team)'' after the name.

\step \act{Select the team\m Manage Certificates...\m +\m Apple Development.}
      \see\ a new line ``Apple Development'' in the list.

\step \act{Read the Team ID and write it down:}
\begin{cmd}
security find-certificate -a -c "Apple Development" -p | openssl x509 -noout -subject
\end{cmd}
      \see\ a line that contains \code{OU=} followed by ten characters. Those ten characters are the Team ID.
\end{jsteps}
\begin{trap}[Free Personal Team rules.]
One signing owner and one bundle id (\code{com.winstonthov.atlas}) for the whole team. A build signed by another account
cannot install over the existing one. An installed build stops launching after 7 days; reinstalling renews it. Ten new
app ids per 7 days. At most three registered devices.
\end{trap}

% =====================================================================
\Job{B}{Prepare the phones}{The iPhone 15 Pro Max first, then the iPhone 14 Pro Max. Do every step for each phone.}
\begin{jsteps}
\step \act{Plug the phone into the Mac, unlock it, tap Trust.}

\step \act{In Xcode: Window\m Devices and Simulators.} Select the phone in the list.

\step \act{On the phone: Settings\m Privacy \& Security\m Developer Mode\m on.} The phone restarts; unlock it and confirm.

\step \act{Back in Xcode, wait until it has finished preparing the phone.} Then check:
\begin{cmd}
xcrun devicectl list devices
\end{cmd}
      \see\ the phone with the state \code{available (paired)}. \code{unavailable} means it is not plugged in.
      \code{connected (no DDI)} usually means Developer Mode is off or Xcode has not finished preparing it: repeat B2 and B3.

\step \act{Write down the phone's iOS version} (Settings\m General\m About) for \code{docs/versions.md}.

\step \act{Put the phone in the headset:} the rear camera and the LiDAR must not be covered.
\end{jsteps}
\begin{later}[Note.]
The Mac also remembers an iPhone 16e. It has no LiDAR and is not one of the plan's phones.
\end{later}

% =====================================================================
\Job{C}{First build to the phones}{Do it once A and B are done and the lanes A and C session reports that the project compiles.}
\begin{jsteps}
\step \act{Close the project in the Unity Editor.}

\step \act{In Terminal, from the repository root.} Replace the X's with the Team ID; it is needed on the first build only:
\begin{cmd}
ATLAS_TEAM_ID=XXXXXXXXXX tools/build_ios.sh
\end{cmd}
      \see\ the app installed on every connected phone. Logs are in \code{ATLAS/Build/}.

\step \act{If the script cannot install:} open \code{ATLAS/Build/iOS/Unity-iPhone.xcodeproj} in Xcode and press Run.

\step \act{First launch on each phone,} with a free account: Settings\m General\m VPN \& Device Management\m trust the
      developer profile.

\step \act{Accept the camera prompt} with the phone in your hand, not in the headset.
\end{jsteps}
\begin{later}[Note.]
The build script and the app code are still uncommitted work of that session. A compile error there is its code
mid-change, not your signing.
\end{later}

% =====================================================================
\Job{D}{Decisions only you can make}{Four decisions. None is urgent this hour.}
\begin{jsteps}
\step \act{The organisers' rules.} Needed before the contracts freeze (gate G0).
\begin{ticks}
\item Read the event's rules and FAQ, and the challenge brief in the repository
      (\cc!documents/Huawei_Custom_Challenge_1_Beyond_Euclid_2016.docx.pdf!), for what they say about work done before the event.
\item If they do not say, ask the organisers in writing: may code and assets be prepared before the event; what must be
      declared as prior work; how long a slot is; one judge or a panel; live demo or video; the space, light and power at
      the venue; whether our own Wi-Fi router is allowed.
\item Tell the team the answer and write it under open item 15 in section 14 of the master plan.
\end{ticks}

\step \act{The headset.} Needed by gate G4.
\begin{ticks}
\item Find the model name, on the headset or its box.
\item Get its three lens numbers: the distortion coefficients, the lens separation and the screen-to-lens distance.
\item Check the fit with both phones and with someone who wears glasses.
\item Give the model and the numbers to lane A. The lens correction (task \textbf{P1-A05}) needs them; without them
      straight lines bow.
\end{ticks}

\step \act{A hosted image-to-3D service,} for when TripoSG fails on the Station. Needed by gate G6.
\begin{ticks}
\item Start from the result of the TripoSG test on the Station (Part 6 of the Station guide,
      \cc!documents/ATLAS_Station_Setup.pdf!): whether TripoSG runs on the RTX 5070, and how many seconds one object takes.
\item Choose the service. For each candidate note its cost, the licence of what it returns, and how long it keeps the
      photos sent to it.
\item Create the account and the key, and give the key to lane D. Never commit it.
\item Tell lane D what a walker must be told before the scan: with a hosted service, crops of the object photos leave the
      Station.
\end{ticks}

\step \act{The non-commercial licences inside the TripoSG stack.} Needed before the Forge's package list is frozen.

      \textbf{The facts:} TripoSG is MIT. \code{diso} is CC BY-NC 4.0. TripoSG's own background remover (RMBG-1.4) is
      non-commercial. The default model of \code{rembg} (\code{bria-rmbg}) is non-commercial too.
\begin{ticks}
\item Check whether the event or the challenge restricts non-commercial parts.
\item Decide: accept them for the entry, or have lane D avoid them.
\item Either way, have lane D enter all four in the licence ledger.
\end{ticks}
\end{jsteps}

% =====================================================================
\Job{E}{The agent sessions}{Nothing else to do by hand.}
\noindent Two sessions are writing code: lanes A and C on the branch \code{phase1-acd}, and the lung data pipeline on the
branch \code{p2-lung-data}. At 19:50 neither branch was on GitHub: \code{phase1-acd} had one commit that was not pushed and 173 uncommitted files, and \code{p2-lung-data} had no commit.\par
\vspace{0.5mm}
\begin{jsteps}
\step \act{In each of the two sessions, paste this line:}
\begin{cmd}
Commit your work to your branch and push the branch to origin. Do not push to main.
\end{cmd}

\step \act{Check that both branches reached GitHub.} In Terminal, from the repository root:
\begin{cmd}
git fetch origin
git branch -r
\end{cmd}
      \see\ \code{origin/phase1-acd} and \code{origin/p2-lung-data} in the list.

\step \act{When a session reports that it has finished,} read its report and have its branch reviewed before it is merged
      into main.
\end{jsteps}

\end{document}
